\RequirePackage{pdfmanagement}
\SetKeys[document/metadata]{pdfversion=1.7, pdfstandard=A-2u, lang=en-GB}
\input{preamble/pdfA-compliance/glyphtounicode}
\documentclass[11pt,a4paper]{report}

\newcommand*{\PaperTitle}{Four limit cycles in a planar three-reaction mass-action system: a proof of the Boros--Hofbauer conjecture}
\newcommand*{\PaperTitlePlain}{Four limit cycles in a planar three-reaction mass-action system: a proof of the Boros-Hofbauer conjecture}
\newcommand*{\PaperShortTitle}{Four limit cycles with three reactions}
\newcommand*{\PaperSubtitle}{chemistry-open-problems \textperiodcentered{} 18 October 2026 \textperiodcentered{} independently verified}
\newcommand*{\PaperCredits}{Solver: Claude Opus 5.5 (effort xhigh), sub-agent of the orchestrator; literature check by a nested Claude Opus 5.5 sub-agent (general-purpose, effort xhigh) \textperiodcentered{} LaTeX edition: Claude Opus 5.5 (effort xhigh) \textperiodcentered{} Independent verification: Claude Opus 5.5 (effort xhigh), fresh context, verdict CONFIRMED WITH CORRECTIONS (minor; see verification.md)}
\newcommand*{\PaperKeywords}{mass-action systems, limit cycles, focal values, Bautin bifurcation, interval arithmetic}
\input{chemop-preamble}

\begin{document}
\makepapertitle

\begin{paperabstract}
Boros and Hofbauer conjectured that the planar mass-action system of three reactions
\(\dot x=1/y-x^ay^b\), \(\dot y=K(-(d-1)-1/y+dx^ay^b)\) has, for suitable parameters, an asymptotically stable
equilibrium surrounded by four limit cycles, after locating numerically a point where the first three focal values
vanish and the fourth is negative; they could not handle \(L_2,L_3,L_4\) analytically. We prove the conjecture by two
independent computer-assisted arguments. First, we compute the focal values exactly as polynomials over
\(\mathbb Q\) in a square-root-free parametrisation, certify with a Krawczyk operator in ball arithmetic a simple
common zero of \(L_1,L_2,L_3\) with \(L_4=-2.580126499494\ldots\cdot10^{-6}<0\) (rigorous enclosure of radius
\(10^{-34}\)), and conclude by successive
perturbations. Second, for explicit decimal parameters, validated Taylor integration of the first-return map brackets
four sign changes of the displacement function, giving four nested limit cycles with alternating stability. We also
show that this order-four weak focus is the only point of the parameter domain where \(L_1=L_2=L_3=0\).
\end{paperabstract}

\section{The problem}

B. Boros, J. Hofbauer, \emph{Oscillations in planar deficiency-one mass-action systems}, J. Dyn. Diff. Equat. 36
(Suppl 1):175--197 (2024), doi:10.1007/s10884-021-10051-z, arXiv:2103.00972 (also PMC10901957), §5.1
(``Four limit cycles''). They consider the three-reaction network (15) with
\((a_1,b_1)=(0,0)\), \((a_2,b_2)=(0,-1)\), \((a_3,b_3)=(a,b)\), \((c_1,d_1)=(0,-1)\), \((c_2,d_2)=(1,-1)\),
\((c_3,d_3)=(-1,d)\), \(a>0\), \(b>-1\), \(d>0\), \(1+bd>0\). After scaling the equilibrium to \((1,1)\)
(possible iff \(d>1\)) the mass-action ODE (16) becomes

\[\dot x = \frac1y - x^a y^b,\qquad
\dot y = K\Big(-(d-1) - \frac1y + d\,x^a y^b\Big),\qquad K>0. \tag{BH}\]

Verbatim (LaTeX source of arXiv:2103.00972, \texttt{sections/5\_three\_reactions.tex}, §5.1; the same text is in the
published version):

\begin{quote}
``In this subsection we discuss why we strongly conjecture that there exist parameter values for which the
differential equation (16) has at least 4 limit cycles. {[}\ldots{]} Set \(K=\frac{a}{1+bd}\) to make \(\operatorname{tr} J\)
equal to zero. The first focal value, \(L_1\), is then
\(L_1 = \frac{\pi a (2a^2d+a(1+bd)-(1+bd)^2)}{8a\sqrt{(d-1)(1+bd)^3}}\) {[}\ldots{]} Thus, there exist \(a\), \(b\), \(d\)
such that \(L_1=L_2=L_3=0\) (namely, we numerically find that this happens at \(a = \widehat{a} \approx 1.01282\),
\(b = \widehat{b} \approx 0.65463\), \(d = \widehat{d} \approx 3.28862\)). Since, again numerically, we see that
\(L_4\) is negative for \(\widehat{a}\), \(\widehat{b}\), \(\widehat{d}\), we conjecture that there exist parameter
values for which the unique positive equilibrium of the differential equation (16) is asymptotically stable,
and is surrounded by four limit cycles \(\Gamma_0\), \(\Gamma_1\), \(\Gamma_2\), \(\Gamma_3\), which are unstable,
stable, unstable, stable, respectively. Since the formulas for \(L_2\), \(L_3\), and \(L_4\) get complicated, we
cannot handle them analytically. This is why we leave the existence of four limit cycles a conjecture.''
\end{quote}

\section{Results}

Parameter domain: \(\mathcal P=\{p=(K,a,b,d):\ K>0,\ a>0,\ b>-1,\ d>1,\ 1+bd>0\}\) (open).
For \(p\in\mathcal P\), \((1,1)\) is the unique positive equilibrium of (BH), with Jacobian

\[J(p)=\begin{pmatrix}-a & -(1+b)\\ Kda & K(1+bd)\end{pmatrix},\qquad
\det J = Ka(d-1)>0,\qquad \operatorname{tr}J = -a+K(1+bd).\]

\(L_k\) denotes the \(k\)-th focal value (Lyapunov quantity) at \((1,1)\) on the hypersurface
\(\mathcal T_0=\{\operatorname{tr}J=0\}=\{K=a/(1+bd)\}\), in the normalisation of Boros--Hofbauer (§4.4 below
shows that our \(L_k\) coincide with theirs on all formulas and values they give).

\textbf{Theorem A (four limit cycles).} There is a parameter point
\(p^*=(K^*,a^*,b^*,d^*)\in\mathcal T_0\) with

\[\begin{aligned}
a^* &= 1.0128219986803146000682082688990\ldots, &
b^* &= 0.6546256211042382107272590334359\ldots,\\
d^* &= 3.2886104393099357877241613562326\ldots, &
K^* &= a^*/(1+b^*d^*) = 0.3212443603989608096064135635919\ldots
\end{aligned}\]

(truncations of Arb enclosures whose radii are \(\le 2\cdot10^{-39}\), see
\texttt{certificates/certificate\_output.json}) such that

\begin{enumerate}
\def\labelenumi{\arabic{enumi}.}
\tightlist
\item
  \(L_1(p^*)=L_2(p^*)=L_3(p^*)=0\);
\item
  \(L_4(p^*)\in[-2.5801264994944244863565347627\cdot 10^{-6}\pm 10^{-34}]\), in particular \(L_4(p^*)<0\);
\item
  the Jacobian matrix \(\partial(L_1,L_2,L_3)/\partial(a,b,d)\) at \((a^*,b^*,d^*)\) is nonsingular.
\end{enumerate}

Consequently, every neighbourhood of \(p^*\) in \(\mathcal P\) contains parameters \(p\) for which the
equilibrium \((1,1)\) of (BH) is a hyperbolic, asymptotically stable focus surrounded by (at least) four limit
cycles \(\Gamma_0,\Gamma_1,\Gamma_2,\Gamma_3\) (nested, \(\Gamma_0\) innermost), where \(\Gamma_0\) and \(\Gamma_2\)
are repelling (unstable) and \(\Gamma_1\), \(\Gamma_3\) are attracting (stable). The four cycles can be taken
inside any prescribed neighbourhood of \((1,1)\).

\textbf{Corollary (the conjecture of Boros--Hofbauer).} For the network of §5.1 with parameters \((K,a,b,d)\) as
produced by Theorem A (or the explicit \(p^\sharp\) of Theorem C below) and rate constants \(\kappa_1=d-1\),
\(\kappa_2=1\), \(\kappa_3=K^{-a}\), the mass-action system (16),

\[\dot x = \kappa_2 y^{-1}-\kappa_3x^ay^b,\qquad \dot y=-\kappa_1-\kappa_2y^{-1}+d\,\kappa_3x^ay^b,\]

has the unique positive equilibrium \((K,1)\), which is asymptotically stable and surrounded by four limit
cycles that are unstable, stable, unstable, stable (from the inside out). Indeed, \(x=KX\), \(y=Y\),
\(t=K\tau\) transforms it into (BH) (a linear change of coordinates and a constant positive time rescaling).
Hence the conjecture of §5.1 holds, in the precise form stated there.

\textbf{Theorem B (uniqueness; bonus).} In the whole domain \(\mathcal P\), the set
\(\{\operatorname{tr}J=0,\ L_1=L_2=L_3=0\}\) is the single point \(p^*\). Consequently, for every
\(p\in\mathcal T_0\) the equilibrium \((1,1)\) is a weak focus of order at most 4 (order exactly 4 only at
\(p^*\)); in particular it is never a centre in this family.

\textbf{Theorem C (explicit parameters; second, independent proof).} Let \(p^\sharp=(K,a,b,d)\) be the exact
decimal fractions

\[\begin{aligned}
a&=1.01282282756688503867785684703002815917702655531627846629837,\\
b&=0.654626809470286327709996401918672391083172153706469077547131,\\
d&=3.28860546617690429718388741565996541608656487798103447854874,\\
K&=0.321244556816175423011532938999913745383501015965226302231941
\end{aligned}\]

(\texttt{certificates/explicit\_params.json}; \(\operatorname{tr}J=-1.90\cdot10^{-22}<0\)). Then the equilibrium \((1,1)\)
of (BH) is a hyperbolic stable focus and the system has (at least) four limit cycles
\(\Gamma_0,\dots,\Gamma_3\), crossing \(\Sigma=\{(1+\rho,1):\rho>0\}\) at \(\rho\in(0.002,0.004)\),
\((0.004,0.006)\), \((0.006,0.012)\), \((0.012,0.024)\) respectively, \(\Gamma_0,\Gamma_2\) repelling and
\(\Gamma_1,\Gamma_3\) attracting. The proof (§6.4) uses only validated (interval) Taylor integration of the ODE
and planar topology; it does not use the focal values.

\emph{Remark (not used, not re-proved).} With the classical finite-cyclicity theory of a weak focus (Bautin 1952;
see e.g.~Roussarie 1998, Han--Yu 2012), Theorem B implies that at most four small-amplitude limit cycles can
bifurcate from \((1,1)\) in this family, and that for parameters close enough to \(p^*\) as in the proof of
Theorem A there are exactly four such cycles, all hyperbolic. Our proofs of Theorems A and B do not depend on
this remark.

What is proved by what:

{\def\LTcaptype{none} % do not increment counter
\begin{longtable}[]{@{}
  >{\raggedright\arraybackslash}p{(\linewidth - 4\tabcolsep) * \real{0.3333}}
  >{\raggedright\arraybackslash}p{(\linewidth - 4\tabcolsep) * \real{0.3333}}
  >{\raggedright\arraybackslash}p{(\linewidth - 4\tabcolsep) * \real{0.3333}}@{}}
\toprule\noalign{}
\begin{minipage}[b]{\linewidth}\raggedright
claim
\end{minipage} & \begin{minipage}[b]{\linewidth}\raggedright
method
\end{minipage} & \begin{minipage}[b]{\linewidth}\raggedright
status
\end{minipage} \\
\midrule\noalign{}
\endhead
\bottomrule\noalign{}
\endlastfoot
exact polynomials \(G_k\) (\(k\le4\)) defining \(L_k\) & exact rational arithmetic (FLINT) & proved, modulo correctness of a \textasciitilde250-line program (validated in §4.4, §8; recomputed exactly by the independent verifier with a different algorithm) \\
\(L_1=L_2=L_3=0\) at a unique point of a box, Jacobian nonsingular & Krawczyk operator in Arb ball arithmetic & proved (computer-assisted) \\
\(L_4<0\) there & Arb enclosure & proved (computer-assisted) \\
sign lemma, return map, successive perturbations & pen-and-paper (§4.2, §6) & proved \\
Theorem B & exact resultants + Arb root enclosures & proved (computer-assisted) \\
Theorem C (explicit rational parameters, four cycles bracketed on \(\Sigma\)) & validated Taylor integration in Arb + Jordan-curve/Poincaré--Bendixson argument & proved (computer-assisted), independent of the focal-value code \\
location of the cycles to more digits (\(\rho\approx0.004,0.008,0.016,0.032\) on another ray) & 60-digit floating-point Taylor integration & numerical illustration only \\
\end{longtable}
}

\section{Set-up and reduction}

\textbf{3.1 Sqrt-free parametrisation of \(\mathcal T_0\).} On \(\mathcal T_0\) the eigenvalues of \(J\) are
\(\pm i\omega\) with \(\omega^2=\det J=a^2(d-1)/(1+bd)\). Put

\[c:=1+b>0,\qquad s:=\frac{\omega^2}{a^2}=\frac{d-1}{1+bd}>0,\qquad w:=\sqrt s,\qquad \rho:=1-(c-1)s .\]

Then \(\rho=(1+b)/(1+bd)>0\), \(d=(1+s)/\rho\), \(1+bd=c/\rho\), \(K=a\rho/c\), \(K(d-1)=as\),
\(Kd=a(1+s)/c\). The map \(\Psi:(a,c,s)\mapsto (K,a,b,d)=(a\rho/c,\,a,\,c-1,\,(1+s)/\rho)\) is a real-analytic
bijection from \(\mathcal D=\{a>0,\ c>0,\ s>0,\ \rho>0\}\) onto \(\mathcal T_0\); the map
\((a,b,d)\mapsto(a,c,s)\) has Jacobian determinant \(\partial s/\partial d=(1+b)/(1+bd)^2>0\).

\textbf{3.2 Normal coordinates.} Let \(u=x-1\), \(v=y-1\), \(I(v)=(1+v)^{-1}\), \(E(u,v)=(1+u)^a(1+v)^b\),
\(f=I-E\). Using \(\dot y=K[(d-1)(I-1)-d f]\), the linear change and time rescaling

\[u=cX,\qquad v=a(wY-X),\qquad \tau=\omega t=a w\,t\]

transform (BH) (for \(p=\Psi(a,c,s)\)) into

\[X'=-Y+P(X,Y),\qquad Y'=X+Q(X,Y),\qquad
P=\frac{N_f}{acw},\quad Q=\frac{cN_I-N_f}{ac}, \tag{N}\]

where \(N_I,N_f\) are the parts of degree \(\ge2\) in \((X,Y)\) of \(I\) and \(f\) (direct check: the linear
parts are \(-Y\) and \(X\)). Every Taylor coefficient of \(N_I,N_f\) is a polynomial in \((a,c,w)\) because
\((1+u)^a(1+v)^{c-1}=\sum_{i,j}\binom{a}{i}\binom{c-1}{j}u^iv^j\). The series converge for \(|u|,|v|<1\), so
(N) is real-analytic near the origin, jointly in the parameters.

\section{Focal values}

\textbf{4.1 Algorithm (complex Lyapunov-function method).} Write \(z=X+iY\); (N) reads \(z'=iz+R(z,\bar z)\),
\(R=P+iQ=\sum_{p+q\ge2}R_{pq}z^p\bar z^q\). We look for a real polynomial
\(F=z\bar z+\sum_{3\le p+q\le N}f_{pq}z^p\bar z^q\) (\(f_{qp}=\overline{f_{pq}}\)). Along (N),

\[\dot F=\sum i(p-q)f_{pq}z^p\bar z^q+\sum f_{rs}\big(r z^{r-1}\bar z^{s}R+s z^r\bar z^{s-1}\bar R\big).\]

Let \(S_{pq}\) be the coefficient of \(z^p\bar z^q\) in the second sum computed with the \(f\) of degree
\(<p+q\) (only those contribute, because \(R=O(|z|^2)\)). Degree by degree, \(n=3,4,\dots\):

\[f_{pq}=\frac{i\,S_{pq}}{p-q}\ (p\ne q),\qquad V_k:=S_{k+1,k+1}\in\mathbb R\ (n=2k+2),\qquad f_{k+1,k+1}:=0 .\]

Then, for every \(j\ge1\), the polynomial \(F_j\) built up to degree \(2j+2\) satisfies

\[\dot F_j=\sum_{k=1}^{j}V_k\,|z|^{2k+2}+O(|z|^{2j+3}). \tag{4.1}\]

(The terms of degree \(\le 2j+2\) of \(\dot F_j\) are exactly the ones cancelled or recorded by the recursion;
all remaining terms have degree \(\ge2j+3\) because \(F_j\) is a polynomial and \(R=O(|z|^2)\) is analytic.)
The choice \(f_{k+1,k+1}=0\) is a convention; Lemma 4.2 holds for any choice.

\textbf{4.2 Sign lemma.} \emph{Let \(p\in\mathcal T_0\) and suppose \(V_1(p)=\dots=V_{j-1}(p)=0\ne V_j(p)\). Let \(\Sigma\) be
any ray issuing from \((1,1)\) and transversal to the flow near \((1,1)\) (e.g.~\(\{(1+\rho,1):\rho>0\}\)), and
\(\delta(\rho)=\mathcal P(\rho)-\rho\) the displacement function of the first-return map on \(\Sigma\) (§6.1).
Then \(\operatorname{sign}\delta(\rho)=\operatorname{sign}V_j(p)\) for all sufficiently small \(\rho>0\).}

\emph{Proof.} By (4.1), \(\dot F_j=V_j|z|^{2j+2}(1+O(|z|))\), so \(\dot F_j\) has the sign of \(V_j\) on a punctured
disc \(0<|z|<r_0\). Since \(F_j=|z|^2+O(|z|^3)\), \(F_j\) is positive definite and its restriction to the ray
\(\Sigma\) is strictly increasing in \(\rho\) for small \(\rho\) (\(F_j=q\rho^2+O(\rho^3)\), \(q>0\), in any linear
coordinates). For \(\rho\) small, the orbit from \(\rho\) stays in the punctured disc during one turn (§6.1), so
\(F_j(\mathcal P(\rho))-F_j(\rho)=\int\dot F_j\,dt\) has the sign of \(V_j\), hence so does
\(\mathcal P(\rho)-\rho\). \(\square\)

(When \(V_1=\dots=V_{j-1}=0\), (4.1) gives \(d(|z|^2)/d\tau\approx V_j|z|^{2j+2}\), i.e.
\(\delta(r)=\pi V_j r^{2j+1}+O(r^{2j+2})\) in the coordinate \(r=|z|\) of (N). Using the unscaled coordinate
\(X=u\) instead of \(u=cX\) gives the normalisation of the paper, \(L_j=\pi V_j^{(u)}\); see §4.3.)

\textbf{4.3 Exact polynomials.} Put \(D=acw\) and \(\tilde R:=D\,R=N_f+iw(cN_I-N_f)\) (polynomial coefficients).
Running the recursion on \(\tilde f_{pq}:=D^{\,p+q-2}f_{pq}\) gives the polynomial recursion

\[\tilde S_n=\sum_{m=2}^{n-1}D^{\,n-1-m}\big[\tilde F_m\otimes\tilde R_{n-m+1}\big],\qquad
\tilde f_{pq}=\frac{i\tilde S_{pq}}{p-q},\qquad \tilde V_k=\tilde S_{k+1,k+1}=D^{2k}V_k ,\]

so that \(\tilde V_k\in\mathbb Q[a,c,w]\). The program \texttt{scripts/focal\_exact.py} (python-flint \texttt{fmpq\_mpoly},
complex numbers as pairs of polynomials) computes \(\tilde V_1,\dots,\tilde V_4\) in 0.1 s and checks along the
way that every \(\tilde S_{k+1,k+1}\) is real and that \(\tilde S_{qp}=\overline{\tilde S_{pq}}\). Every
\(\tilde V_k\) is \(w\) times a polynomial in \((a,c,w^2)\); we set

\[G_k(a,c,s):=\tilde V_k(a,c,\sqrt s)/\sqrt s\in\mathbb Q[a,c,s],\qquad
L_k=\pi\,V_k^{(u)}=\frac{\pi\,G_k(a,c,s)}{a^{2k}c^{4k}s^{\,k-1/2}},\]

where \(V_k^{(u)}\) is the Lyapunov quantity in the coordinates \(X=u\) (unscaled), i.e.~\(V^{(u)}_k=c^{-2k}V_k\).
On \(\mathcal D\), \(\operatorname{sign}L_k=\operatorname{sign}V_k=\operatorname{sign}G_k\), so Lemma 4.2 holds with
\(G_k\) in place of \(V_k\). Factorisations (exact):

\[G_1=-\tfrac18 a^2c^2\,g_1,\quad G_2=\tfrac1{1152}a^4c^4\,g_2,\quad G_3=-\tfrac1{442368}a^6c^6\,g_3,\quad
G_4=\tfrac1{530841600}a^8c^8\,g_4,\]

with \(g_k\in\mathbb Z[a,c,s]\) irreducible, of degrees \((2k,2k,3k-1)\) in \((a,c,s)\) and 9, 48, 136, 296
terms. Explicitly

\[-g_1=2a^2(1+s)\rho+ac\rho-c^2,\qquad \rho=1+s-cs,\]

and \(g_2,g_3\) are listed in Appendix A; all four are stored in \texttt{certificates/G\_k.json}.

\textbf{4.4 Validation against the paper.} (i) Exact identity (\texttt{crosscheck.py}): \(L_1=\pi G_1/(a^2c^4\sqrt s)
=\pi[2a^2(1+s)\rho+ac\rho-c^2]/(8c^2\sqrt s)\) equals the paper's
\(\pi a(2a^2d+a(1+bd)-(1+bd)^2)/(8a\sqrt{(d-1)(1+bd)^3})\) after \(d=(1+s)/\rho\), \(1+bd=c/\rho\).
(ii) The generic floating-point implementation \texttt{focal\_mp.py} (independent code path: numerical Taylor
expansion of arbitrary generalized-monomial systems, numerical normalisation with \(\omega=\sqrt{\det J}\), no
reparametrisation) reproduces \textbf{exactly} (ratio \(1.000\ldots\) to 15+ digits) every focal value printed in
the paper: \(L_1\) of §3.2, §4.1 and §5.1 (two generic points each), \(L_2\) of §4.1 along \(\{L_1=0\}\) (two
points), \(L_2\approx0.01293\) of §3.2, and \(L_3=-\frac{625\pi}{110592}\sqrt{7/2}\) of §4.1. (iii) The exact
\(L_k\) agree with \texttt{focal\_mp.py} at four parameter points to relative precision \(10^{-55}\) for \(k=1..4\).
(iv) The paper's checkpoints: at \((a,d)=(1,165/49)\) (then \(b=7/11\), \(s=58/77\)) we get exactly
\(L_1=L_2=0\) and \(L_3<0\) (exact rational sign); at \((a,d)=((1+\sqrt{3961})/60,3)\), \(L_1=L_2=0\) (to 50
digits) and \(L_3>0\) --- as stated in the paper.

\section{\texorpdfstring{The interval certificate (\texttt{scripts/certify.py})}{The interval certificate (scripts/certify.py)}}

\textbf{Krawczyk's theorem} (Krawczyk 1969; Moore 1977; Neumaier 1990, §5.1). Let \(G:\mathbb R^n\to\mathbb R^n\) be
\(C^1\) on a box \(X\), \(m\in X\), \(C\) a real \(n\times n\) matrix and \(\mathbf J\) an interval matrix
containing \(G'(x)\) for all \(x\in X\). If
\(\mathcal K(X)=m-CG(m)+(I-C\mathbf J)(X-m)\subset\operatorname{int}X\), then \(G\) has exactly one zero in
\(X\). Moreover, if \(\|I-C\mathbf J\|_\infty<1\), every matrix in \(\mathbf J\) is nonsingular (Neumann series
for \(CA\)).

Data (\texttt{certificates/box.json}): variables \((a,c,s)\), centre \(m\) = three rationals with denominator
\(\le 10^{60}\) (Newton iterate at 150 digits), radius \(r=10^{-40}\), preconditioner \(C\) = rational
approximation of \(G'(m)^{-1}\). Computations in Arb (python-flint 0.9.0, 512-bit midpoint--radius balls with
outward rounding); \(G=(G_1,G_2,G_3)\), \(\mathbf J\) = exact derivative polynomials evaluated on \(X\).

{\def\LTcaptype{none} % do not increment counter
\begin{longtable}[]{@{}
  >{\raggedright\arraybackslash}p{(\linewidth - 2\tabcolsep) * \real{0.5000}}
  >{\raggedright\arraybackslash}p{(\linewidth - 2\tabcolsep) * \real{0.5000}}@{}}
\toprule\noalign{}
\begin{minipage}[b]{\linewidth}\raggedright
check
\end{minipage} & \begin{minipage}[b]{\linewidth}\raggedright
result
\end{minipage} \\
\midrule\noalign{}
\endhead
\bottomrule\noalign{}
\endlastfoot
C0 recomputed \(G_k\) = stored \(G_k\); C0′ Arb balls contain the exact box & PASS \\
C1 on \(X\): \(a>0\), \(c>0\), \(s>0\), \(\rho>0\), \(d>1\), \(1+bd>0\), \(K>0\) & PASS \\
C2 \(\mathcal K(X)\subset\operatorname{int}X\): \(\max_i|\mathcal K_i-m_i|\le3.62\cdot10^{-61}\ll r=10^{-40}\) & PASS \\
C3 \(\|I-C\mathbf J\|_\infty\le1.37\cdot10^{-32}<1\); \(\det\mathbf J\in[-0.208483123515217\pm4\cdot10^{-16}]\) & PASS \\
C4 \(G_4(X)\subset[-9.354314437536414185821767970\cdot10^{-4}\pm2\cdot10^{-32}]<0\); \(L_4(X)\subset[-2.5801264994944244863565347627\cdot10^{-6}\pm10^{-34}]\) & PASS \\
C5 \(\det\partial(G_1,G_2,G_3)/\partial(a,b,d)\in[-0.03470367143\pm4\cdot10^{-12}]\) & PASS \\
\end{longtable}
}

Let \(\mu^*\in X\) be the unique zero and \(p^*=\Psi(\mu^*)\). At \(\mu^*\), \(G_1=G_2=G_3=0\), so
\(L_1=L_2=L_3=0\); \(L_4(p^*)<0\) by C4; and since \(L_k=\pi G_k/(\text{positive})\), at a common zero
\(\partial L_k=\pi\,\partial G_k/(\text{positive})\), so C5 gives item 3 of Theorem A. The enclosures of
\((a^*,b^*,d^*,K^*)\) quoted in §2 are the Arb images of \(X\) (\texttt{certificates/certificate\_output.json}). The
paper's values \(a\approx1.01282\), \(b\approx0.65463\), \(d\approx3.28862\) agree with ours except in the
last digit of \(d\) (\(d^*=3.288610\ldots\)), which reflects their numerical accuracy.

\section{The bifurcation argument (proof of Theorem A) and validated integration (proof of Theorem C)}

Let \(\Sigma=\{(1+\rho,1):\rho>0\}\). On \(\Sigma\), \(\dot y=Kd((1+\rho)^a-1)>0\): \(\Sigma\) is transversal.

\textbf{6.1 Lemma (return map).} Let \(p_0\in\mathcal T_0\). There are a compact neighbourhood \(U\subset\mathcal P\) of
\(p_0\) and \(\bar\rho>0\) such that for all \(p\in U\) and \(\rho\in(0,\bar\rho]\) the forward orbit of
\((1+\rho,1)\) returns to \(\Sigma\) after one turn around \((1,1)\), at \((1+\mathcal P(\rho,p),1)\); the map
\(\mathcal P\) is continuous on \((0,\bar\rho]\times U\), real-analytic in \(\rho\), and strictly increasing in
\(\rho\). Each zero \(\rho_0\in(0,\bar\rho]\) of \(\delta(\cdot,p)=\mathcal P(\cdot,p)-\mathrm{id}\) gives a
periodic orbit through \((1+\rho_0,1)\) surrounding \((1,1)\); distinct zeros give distinct periodic orbits, and an
isolated zero gives an isolated periodic orbit, i.e.~a limit cycle.

\emph{Proof (standard).} Let \(T\) bring \(J(p_0)\) to the rotation form and use polar coordinates \((r,\theta)\) for
\(\xi=T(u,v)\). Since \(J(p)\to J(p_0)\) as \(p\to p_0\) and the nonlinear terms are \(O(r^2)\) uniformly on a
compact parameter set, there are \(U\), \(r_0>0\) with \(\dot\theta\ge\omega_0/2\) and \(|\dot r|\le Lr\) for
\(r\le r_0\), \(p\in U\). For \(r(0)\le\bar r:=r_0e^{-4\pi L/\omega_0}\) the orbit stays in \(r\le r_0\) up to
time \(4\pi/\omega_0\), during which \(\theta\) increases by at least \(2\pi\); hence it meets the ray \(\Sigma\)
(a ray in the \(\xi\)-plane too) again. The crossings are transversal, so the implicit function theorem gives
continuity in \((\rho,p)\) and analyticity in \(\rho\). Monotonicity follows from uniqueness of solutions
(orbits do not cross). A periodic orbit in \(\{r\le r_0\}\) has winding number \(\pm1\) around \((1,1)\) and
\(\theta\) is strictly increasing along it, so it meets \(\Sigma\) exactly once per period; a periodic orbit near
one through \(\rho_0\) meets \(\Sigma\) near \(\rho_0\), hence is the same orbit if \(\rho_0\) is an isolated zero.
\(\square\)

\textbf{6.2 Lemma (trace).} If \(p\in U\) and \(\operatorname{tr}J(p)<0\), then \(\delta(\rho,p)<0\) for all small
\(\rho>0\). \emph{Proof.} \((1,1)\) is a hyperbolic sink (\(\det J>0\)), attracting a neighbourhood. If
\(\delta(\rho,p)=0\) there would be a periodic orbit in that neighbourhood; if \(\delta(\rho,p)>0\), the
iterates \(\mathcal P^n(\rho)\) would increase (monotonicity) and not tend to 0. Both contradict attraction.
\(\square\)

\textbf{6.3 Proof of Theorem A.} Let \(G=(G_1,G_2,G_3)\). By C2--C3, \(G(\mu^*)=0\) and \(G'(\mu^*)\) is invertible,
so (inverse function theorem) \(G\) maps a neighbourhood \(N\) of \(\mu^*\) diffeomorphically onto a
neighbourhood \(W\) of \(0\); write \(\Phi=\Psi\circ(G|_N)^{-1}:W\to\mathcal T_0\), continuous with
\(\Phi(0)=p^*\). Fix \(U,\bar\rho\) from Lemma 6.1 for \(p_0=p^*\) and shrink \(W\) so that \(\Phi(W)\subset U\).
Given any neighbourhood \(\mathcal V\) of \(p^*\) and \(\varepsilon>0\), we construct \(p_4\in\mathcal V\) and
\(0<\rho_5<\rho_4<\rho_3<\rho_2<\rho_1<\varepsilon\) with

\[\delta(\rho_5,p_4)<0<\delta(\rho_4,p_4),\quad \delta(\rho_3,p_4)<0<\delta(\rho_2,p_4),\quad \delta(\rho_1,p_4)<0 .\]

\begin{itemize}
\tightlist
\item
  Step 0. At \(p^*\): \(G_1=G_2=G_3=0\), \(G_4<0\). Lemma 4.2 (\(j=4\)) gives \(\rho_1\in(0,\min(\bar\rho,\varepsilon))\)
  with \(\delta(\rho_1,p^*)<0\).
\item
  Step 1. By continuity of \(\delta(\rho_1,\cdot)\) choose \(t_3>0\) so small that \(p_1:=\Phi(0,0,t_3)\in\mathcal V\)
  still has \(\delta(\rho_1,p_1)<0\). At \(p_1\): \(G_1=G_2=0<G_3\), so by Lemma 4.2 there is \(\rho_2<\rho_1\)
  with \(\delta(\rho_2,p_1)>0\).
\item
  Step 2. Choose \(t_2>0\) so small that \(p_2:=\Phi(0,-t_2,t_3)\in\mathcal V\) keeps the signs at \(\rho_1,\rho_2\).
  At \(p_2\): \(G_1=0>G_2\): choose \(\rho_3<\rho_2\) with \(\delta(\rho_3,p_2)<0\).
\item
  Step 3. Choose \(t_1>0\) so small that \(p_3:=\Phi(t_1,-t_2,t_3)\in\mathcal V\) keeps the signs at
  \(\rho_1,\rho_2,\rho_3\). At \(p_3\): \(G_1>0\): choose \(\rho_4<\rho_3\) with \(\delta(\rho_4,p_3)>0\).
\item
  Step 4. Let \(p_4=p_3-(\eta,0,0,0)\) (decrease \(K\)), \(\eta>0\) so small that \(p_4\in\mathcal V\cap U\) keeps the
  signs at \(\rho_1,\dots,\rho_4\). Now \(\operatorname{tr}J(p_4)=-\eta(1+bd)<0\); Lemma 6.2 gives \(\rho_5<\rho_4\) with
  \(\delta(\rho_5,p_4)<0\).
\end{itemize}

(All the \(\Phi(\cdot)\) points lie on \(\mathcal T_0\), where Lemma 4.2 applies; continuity of \(\Phi\) makes each
new point as close to the previous one as needed.) Since \(\delta(\cdot,p_4)\) is analytic on \((0,\bar\rho]\) and
not identically zero, it has finitely many zeros in \([\rho_5,\rho_1]\). On each of the four intervals
\((\rho_5,\rho_4),(\rho_4,\rho_3),(\rho_3,\rho_2),(\rho_2,\rho_1)\) it changes sign; take the zeros
\(\zeta_0<\zeta_1<\zeta_2<\zeta_3\), one in each interval, at which the sign changes from \(-\) to \(+\)
(\(\zeta_0,\zeta_2\)) resp. from \(+\) to \(-\) (\(\zeta_1,\zeta_3\)) when \(\rho\) increases (they exist because the
finitely many gaps between consecutive zeros carry constant signs that start with \(-\), resp. \(+\), and end
with the opposite sign). Let \(\Gamma_i\) be the periodic orbit through \((1+\zeta_i,1)\) --- a limit cycle by Lemma 6.1.
If \(\delta<0\) just inside and \(\delta>0\) just outside \(\zeta_0\), orbits near \(\Gamma_0\) move away from it on
both sides (\(\mathcal P\) is increasing and fixes \(\zeta_0\)): \(\Gamma_0\) is repelling; likewise \(\Gamma_2\).
At \(\zeta_1,\zeta_3\) the signs are reversed and the iterates converge monotonically to \(\zeta_i\) from both sides:
\(\Gamma_1,\Gamma_3\) are attracting. The equilibrium is a hyperbolic sink (\(\operatorname{tr}J<0<\det J\)). The
\(\Gamma_i\) are nested Jordan curves around \((1,1)\) (each meets \(\Sigma\) once, at increasing \(\zeta_i\)) and lie
in any prescribed neighbourhood of \((1,1)\) if \(\varepsilon\) is small. \(\square\)

\textbf{6.4 Proof of Theorem C (\texttt{scripts/validated\_cycles.py}).}

\emph{Validated integrator.} With \(E=x^ay^b\), \(I=1/y\), \(J=1/x\), (BH) lifts to the polynomial system
\(x'=I-E\), \(y'=K(-(d-1)-I+dE)\), \(E'=E(aJx'+bIy')\), \(I'=-I^2y'\), \(J'=-J^2x'\); from consistent initial data
its solution projects onto that of (BH) (indeed \(q=Iy\) and \(p=Jx\) satisfy \(q'=Iy'(1-q)\), \(p'=Jx'(1-p)\), so
\(q\equiv p\equiv1\) and \(y,x\) stay positive; then \(E'=E(ax'/x+by'/y)\) gives \(E\equiv x^ay^b\)). One step of length
\(h\) from a ball enclosure \(Z_k\) of the state:
(i) find a box \(B\) with \(Z_k+[0,h]\,F(B)\subset\operatorname{int}B\) (then every solution starting in \(Z_k\)
exists on \([0,h]\) and stays in \(B\): if it left \(B\) at a first time \(t_1\le h\), then
\(z(t_1)\in z_0+t_1\operatorname{conv}F(B)\subset\operatorname{int}B\), a contradiction); (ii) Taylor
coefficients \(z_j(Z_k)\), \(j<N\), by the standard recurrences in Arb; (iii) Lagrange remainder
\(z_N(B)\tau^N\) (the \(N\)-th Taylor coefficient is a polynomial function of the state, and the intermediate
state lies in \(B\)). Hence \(Z(t_k+\tau)\in\sum_{j<N}z_j(Z_k)\tau^j+z_N(B)\tau^N\) for \(\tau\in[0,h]\). Parameters:
\(N=36\), \(h=0.05\) (146 steps per turn), 320-bit balls. The enclosures of \(\delta\) do not change (15 digits)
with \((N,h)=(28,0.03)\) or \((44,0.08)\).

\emph{Crossing bookkeeping.} \(\Sigma\) is transversal (\(\dot y=Kd((1+\rho)^a-1)>0\)). In every step whose box \(B\)
meets the line \(y=1\), the script requires either (\(x<1\) and \(\dot y<0\) on \(B\)) or (\(x>1\) and
\(\dot y>0\) on \(B\)), so \(y\) is strictly monotone there and crossings are counted exactly: for every test
point the orbit leaves \(\Sigma\) upwards, crosses the left half-line \(\{x<1,y=1\}\) exactly once (downwards),
and the first return to \(\Sigma\) is enclosed by bisection of the crossing time on the Taylor model
(\(2^{-140}\) relative width), giving \(\mathcal P(\rho)\) to about \(10^{-38}\).

\emph{Certified values} (\texttt{certificates/validated\_cycles\_output.json}):

{\def\LTcaptype{none} % do not increment counter
\begin{longtable}[]{@{}ll@{}}
\toprule\noalign{}
\(\rho\) & \(\delta(\rho)=\mathcal P(\rho)-\rho\) \\
\midrule\noalign{}
\endhead
\bottomrule\noalign{}
\endlastfoot
0.002 & \([-4.60438730999\cdot10^{-25}\pm 6.2\cdot10^{-38}]\) \\
0.004 & \([+1.27356613600\cdot10^{-24}\pm 1.9\cdot10^{-36}]\) \\
0.006 & \([-3.54914581811\cdot10^{-24}\pm 4.5\cdot10^{-36}]\) \\
0.012 & \([+1.41126309912\cdot10^{-22}\pm 4.5\cdot10^{-34}]\) \\
0.024 & \([-3.02911287653\cdot10^{-20}\pm 4.8\cdot10^{-32}]\) \\
\end{longtable}
}

(radii as printed by Arb with 12 significant digits, i.e.~including the display rounding).

(also \(\operatorname{tr}J<0<\det J\), \((\operatorname{tr}J)^2<4\det J\) in exact rational arithmetic).

\emph{Topology.} Let \(\rho_1=0.024\), \(q_1=(1+\rho_1,1)\), \(\gamma\) the orbit arc from \(q_1\) to its first return
\(q_1'=(1+\mathcal P(\rho_1),1)\), \(\mathcal P(\rho_1)<\rho_1\). Use polar coordinates \((r,\theta)\) around
\((1,1)\) in the plane \((u,v)=(x-1,y-1)\), so that \(\Sigma=\{\theta=0\}\). The script certifies
\(u\dot v-v\dot u>0\) on every enclosure box along \(\gamma\), i.e.~\(\theta\) is strictly increasing along
\(\gamma\); with the crossing bookkeeping (\(\gamma\) meets \(\{\theta=\pi\}\) once and \(\{\theta=0\}\) only at its
end points), \(\theta\) increases by exactly \(2\pi\), so \(\gamma\) is a polar graph \(r=R(\theta)\),
\(0\le\theta\le2\pi\), with \(R(0)=\rho_1>R(2\pi)=\mathcal P(\rho_1)\). Hence \(C=\gamma\cup[q_1',q_1]\) is a Jordan
curve bounding the star-shaped region \(D=\{r\le R(\theta),\ 0<\theta<2\pi\}\cup\{\theta=0,\ r\le\rho_1\}\),
which contains \((1,1)\) and the whole segment \(\{(1+\rho,1):0<\rho\le\rho_1\}\). A point of the open segment
\((q_1',q_1)\) has \(R(2\pi^-)=\mathcal P(\rho_1)<r<\rho_1=R(0^+)\): the flow, which crosses \(\Sigma\) with
increasing \(\theta\), enters \(D\) there; orbits cannot cross \(\gamma\). So \(D\) is positively invariant. For every
\(\rho\in(0,\rho_1]\) the orbit of \((1+\rho,1)\) stays in \(D\), where \((1,1)\) is the only equilibrium (it is
the unique positive equilibrium); by Poincaré--Bendixson its \(\omega\)-limit set is \((1,1)\) (a focus: the orbit
spirals and crosses \(\Sigma\)) or a periodic orbit surrounding \((1,1)\) (which meets \(\Sigma\) transversally, so
the orbit crosses \(\Sigma\) near it); a homoclinic loop is impossible since \((1,1)\) is a sink. Thus the first
return map \(\mathcal P\) is defined on \((0,\rho_1]\), and by transversality it is continuous and real-analytic
there. By the table, \(\delta\) changes sign on each of the four intervals; as in §6.3 (finitely many zeros,
choice of sign-changing zeros, monotonicity of \(\mathcal P\)) we get four limit cycles with the stated
stabilities. \(\square\)

The parameters \(p^\sharp\) were found with the focal values (targets
\((G_1,G_2,G_3)\) prescribed so that the truncated displacement has roots at \(0.004,0.008,0.016,0.032\) on the
ray \(Y=0\), see §8.2), but the proof of Theorem C does not depend on how they were found.

\textbf{Remarks.} (a) Only the network parameters \((a,b,d)\) and \(K=\bar x/\bar y\) are varied, exactly as in the
paper's §§4.1 and 5.1; the conjecture is about ``parameter values'' of (16), including the exponents. (For a
fixed network the rate constants enter the scaled system (BH) only through \(K\) --- the scaled constants
\(\overline\kappa_i\) are fixed up to the factor \(\lambda\), which is absorbed by a time rescaling --- so a fixed
network only gives a one-parameter family and the network parameters must be varied, as the paper does; compare
arXiv:2406.13451 for a general statement in the quadratic setting.) (b) The perturbation
directions are the paper's: first \(L_3\), then \(L_2\), \(L_1\), finally the trace via \(K\).

\section{\texorpdfstring{Proof of Theorem B (\texttt{scripts/uniqueness.py})}{Proof of Theorem B (scripts/uniqueness.py)}}

On \(\mathcal D\), \(G_1=0\iff 2a^2(1+s)\rho+ac\rho-c^2=0\). With \(t:=a/c>0\) this reads
\(\rho\,t(2t(1+s)+1)=1\), i.e.~\(\rho=1/(tQ)\), \(Q=2t(1+s)+1\), and \(c=(1+s-\rho)/s\). Hence
\((t,s)\mapsto(a,c,s)=(tc,c,s)\) is a bijection from \(\{t>0,\ s>0,\ 2t(1+s)>1\}\) (the last condition is
\(c>0\)) onto \(\{G_1=0\}\cap\mathcal D\). Substituting and clearing the denominator \((stQ)^m\):

\begin{itemize}
\tightlist
\item
  \(G_1\) becomes identically 0 (check);
\item
  \(H_2=\text{const}\cdot(\text{powers of } s,\ s+1,\ t,\ ts+t+1,\ 2ts+2t-1)\cdot h_2(t,s)\), and the same for
  \(H_3\) with \(h_3\); all the listed small factors are nonzero on the \((t,s)\)-domain
  (\(2t(1+s)-1>0\iff c>0\)); \(h_2\), \(h_3\) have degrees \((6,4)\), \((10,8)\) and 24, 68 terms.
\item
  \(R_t=\operatorname{Res}_s(h_2,h_3)\) (degree 58) factors as
  \(t^{32}(t+1)^9(2t-1)^4(2t+3)(3t-1)\,p_{11}(t)\) with
  \(p_{11}=628224t^{11}+1764800t^{10}-1694592t^9-15308848t^8-25081168t^7+10994160t^6+36659948t^5-2189226t^4-16683601t^3+3894133t^2-2248520t+1436225\);
  \(R_s=\operatorname{Res}_t(h_2,h_3)\) has degree 56.
\item
  Any common zero \((t_0,s_0)\) satisfies \(R_t(t_0)=R_s(s_0)=0\). All complex roots of both resultants are
  enclosed in Arb balls; the candidates with possibly positive real part and possibly zero imaginary part are
  \(t\in\{1/3,\ 0.61212,\ 0.61956,\ 0.74045,\ 2.73230,\ 1/2\}\) and \(s\in\{0.11924,\ 0.72590\}\). For all pairs
  but one, the Arb enclosure of \(h_2\) or \(h_3\) excludes 0. The surviving pair, mapped back to \((a,c,s)\), lies
  inside the Krawczyk box \(X\) of §5, where the zero is unique.
\end{itemize}

Hence \(\{G_1=G_2=G_3=0\}\cap\mathcal D=\{\mu^*\}\), which is Theorem B (the order statement follows from
\(L_4(p^*)\ne0\)). (PARI/GP gives the same picture: over \(\mathbb Q(\theta)\), \(p_{11}(\theta)=0\),
\(\gcd(h_2,h_3)\) is linear in \(s\); the real roots \(\theta\approx0.6196, 0.7405, 2.7323\) give \(s<0\).)

\section{Non-rigorous cross-checks (not part of the proof)}

\begin{enumerate}
\def\labelenumi{\arabic{enumi}.}
\tightlist
\item
  \textbf{Independent method} (\texttt{polar\_check.py}): the coefficients \(u_k(2\pi)\) of the displacement map, obtained
  by integrating the triangular ODE system of the polar expansion \(dr/d\theta=\sum C_k(\theta)r^k\) with DOP853
  (tolerance \(10^{-13}\)), agree with \(\pi V_k\): \(u_3=\pi V_1\) at a generic point (rel. error
  \(6\cdot10^{-15}\)); on \(\{L_1=0\}\): \(u_3\approx0\), \(u_5=\pi V_2\) (\(10^{-11}\)); on \(\{L_1=L_2=0\}\):
  \(u_5\approx0\), \(u_7=\pi V_3\) (\(8\cdot10^{-11}\)); at \(p^*\): \(u_2,\dots,u_8=O(10^{-14})\) and
  \(u_9(2\pi)=-2.5801263\cdot10^{-6}=\pi V_4\) (\(5\cdot10^{-8}\)).
\item
  \textbf{Explicit four cycles} (\texttt{illustrate\_cycles.py}, 60-digit arithmetic, Taylor integrator of order 40 on the
  polynomial lift \((x,y,x^ay^b,1/y,1/x)\)): targets chosen so that the truncated displacement has roots
  \(0.004,0.008,0.016,0.032\); parameters
  \(a=1.012822827566885038677856847030028\ldots\), \(b=0.654626809470286327709996401918672\ldots\),
  \(d=3.288605466176904297183887415659965\ldots\), \(K=a/(1+bd)-6.03\cdot10^{-23}\) (full 60 digits printed by
  the script). Displacements on the ray \(\{x=1+\rho,\ y=1-\rho a/(1+b)\}\):
  \(\delta(0.002)=-9.56\cdot10^{-25}\), \(\delta(0.006)=+1.89\cdot10^{-24}\), \(\delta(0.012)=-3.12\cdot10^{-23}\),
  \(\delta(0.024)=+2.41\cdot10^{-21}\), \(\delta(0.045)=-7.01\cdot10^{-19}\) (truncated predictions within 0.1--15\%);
  bisection locates the four cycles at \(\rho\approx0.004006,\ 0.008019,\ 0.01614,\ 0.03230\). This is floating-point
  evidence only (no validated integration).
\item
  The paper's checkpoints and formulas (§4.4).
\end{enumerate}

\section{Trust base, limitations, what a verifier should check first}

\begin{itemize}
\tightlist
\item
  \textbf{Two independent proofs.} Theorem A rests on the focal-value code + Krawczyk + the bifurcation argument;
  Theorem C rests only on the validated integrator (\textasciitilde150 lines, \texttt{validated\_cycles.py}) + Arb + planar
  topology. Either one alone settles the conjecture. A verifier should check in Theorem C: the a priori
  enclosure test and the remainder term (§6.4), the crossing bookkeeping (no missed crossing of
  \(\Sigma\)), and the Jordan-curve/positive-invariance argument.
\item
  \textbf{Correctness of the focal-value code.} The certificate checks the polynomials \(G_k\) that
  \texttt{focal\_exact.py} produces; it does not prove that these are the Lyapunov quantities. Evidence: the
  derivation in §3--4 (short, checkable by hand), the exact match with the paper's closed-form \(L_1\), the
  match of an independent generic implementation with every focal value printed in the paper (\(L_1\), \(L_2\),
  \(L_3\) of three different systems), and an independent method (polar expansion) at the degenerate points,
  including \(u_9(2\pi)=\pi V_4\) at \(p^*\). A verifier should re-derive \(\tilde R=N_f+iw(cN_I-N_f)\) and, ideally,
  recompute \(L_1,\dots,L_4\) at \(p^*\) with an independent CAS/algorithm (e.g.~Singular/Maple normal form).
  \emph{Done in the independent verification} (\texttt{verification.md}, scripts in \texttt{verification/}): the
  focal values were recomputed exactly with a different algorithm (real-coordinate Lyapunov function, homogeneous
  linear algebra with the circle-average normalisation, independently derived vector field), giving
  \(w\,G_k(a,c,w^2)\) \emph{exactly} for \(k=1,\dots,4\); the Krawczyk certificate was re-done with another interval
  library (mpmath.iv), Theorem B with msolve, and the displacements of Theorem C with an independent validated
  integrator (polar coordinates, angle as time), which reproduces the enclosures of §6.4.
\item
  \textbf{Sign lemma (§4.2) and successive perturbations (§6)}: classical; the proof given is self-contained.
  Check that every perturbed parameter stays on \(\mathcal T_0\) when Lemma 4.2 is used (it does: \(\Phi\) maps into
  \(\mathcal T_0\)), and that the last step only changes \(K\).
\item
  \textbf{Interval arithmetic}: Arb (python-flint 0.9.0); comparisons used are the certain ones (\texttt{\textless{}} true only if
  true for all points of the balls). Polynomials are evaluated term by term (no Horner), which only enlarges
  enclosures.
\item
  \textbf{Mass-action interpretation}: exponents are real (generalized mass-action, real complexes), exactly as in
  the paper's setting (§5: \(a>0\), \(b>-1\), \(d>0\), \(1+bd>0\)).
\item
  \textbf{Not proved}: exact number (=4) of small cycles and their hyperbolicity (classical, cited in the Remark of
  §2; for Theorem C, more cycles could in principle hide in the intervals --- at least four are proved); the
  more precise cycle locations of §8.2 (numerical only); anything about large-amplitude cycles.
\end{itemize}

\section{Reproduction}

From the project root (Git Bash): \texttt{outils/wsl\ "cd\ chemistry-open-problems/problems/boros-hofbauer-four-limit-cycles\ \&\&\ bash\ scripts/run\_all.sh"}.
Individual scripts: \texttt{focal\_exact.py} (exact \(\tilde V_k\)), \texttt{crosscheck.py}, \texttt{certify.py}, \texttt{uniqueness.py},
\texttt{validated\_cycles.py} (each prints PASS/FAIL and exits non-zero on failure; seconds each), \texttt{polar\_check.py}
(\textasciitilde5 s), \texttt{illustrate\_cycles.py} (\textasciitilde3 min, illustration). Full log: \texttt{certificates/run\_all.log} (35 PASS, 0 FAIL).

\section{\texorpdfstring{\(g_1,g_2,g_3\) (\(G_k=\kappa_k a^{2k}c^{2k}g_k\), \(\kappa=(-\frac18,\frac1{1152},-\frac1{442368},\frac1{530841600})\))}{g\_1,g\_2,g\_3 (G\_k=\textbackslash kappa\_k a\^{}\{2k\}c\^{}\{2k\}g\_k, \textbackslash kappa=(-\textbackslash frac18,\textbackslash frac1\{1152\},-\textbackslash frac1\{442368\},\textbackslash frac1\{530841600\}))}}

\texttt{g\_1\ =\ 2*a\^{}2*c*s\^{}2\ +\ 2*a\^{}2*c*s\ +\ a*c\^{}2*s\ -\ 2*a\^{}2*s\^{}2\ -\ 4*a\^{}2*s\ -\ a*c*s\ -\ 2*a\^{}2\ -\ a*c\ +\ c\^{}2}

\texttt{g\_2\ =\ 318*a\^{}4*c\^{}3*s\^{}5\ +\ 396*a\^{}4*c\^{}3*s\^{}4\ +\ 159*a\^{}3*c\^{}4*s\^{}4\ -\ 954*a\^{}4*c\^{}2*s\^{}5\ +\ 78*a\^{}4*c\^{}3*s\^{}3\ +\ 35*a\^{}3*c\^{}4*s\^{}3\ -\ 1998*a\^{}4*c\^{}2*s\^{}4\ -\ 765*a\^{}3*c\^{}3*s\^{}4\ +\ 954*a\^{}4*c*s\^{}5\ -\ 1374*a\^{}4*c\^{}2*s\^{}3\ -\ 1086*a\^{}3*c\^{}3*s\^{}3\ -\ 3*a\^{}2*c\^{}4*s\^{}3\ +\ 3048*a\^{}4*c*s\^{}4\ +\ 1053*a\^{}3*c\^{}2*s\^{}4\ -\ 318*a\^{}4*s\^{}5\ -\ 330*a\^{}4*c\^{}2*s\^{}2\ -\ 409*a\^{}3*c\^{}3*s\^{}2\ -\ 79*a\^{}2*c\^{}4*s\^{}2\ +\ 3804*a\^{}4*c*s\^{}3\ +\ 3391*a\^{}3*c\^{}2*s\^{}3\ +\ 72*a\^{}2*c\^{}3*s\^{}3\ -\ 1446*a\^{}4*s\^{}4\ -\ 447*a\^{}3*c*s\^{}4\ +\ 2280*a\^{}4*c*s\^{}2\ +\ 3511*a\^{}3*c\^{}2*s\^{}2\ +\ 900*a\^{}2*c\^{}3*s\^{}2\ -\ 175*a*c\^{}4*s\^{}2\ -\ 2748*a\^{}4*s\^{}3\ -\ 2340*a\^{}3*c*s\^{}3\ -\ 69*a\^{}2*c\^{}2*s\^{}3\ +\ 570*a\^{}4*c*s\ +\ 1173*a\^{}3*c\^{}2*s\ +\ 636*a\^{}2*c\^{}3*s\ +\ 33*a*c\^{}4*s\ -\ 2748*a\^{}4*s\^{}2\ -\ 4134*a\^{}3*c*s\^{}2\ -\ 743*a\^{}2*c\^{}2*s\^{}2\ +\ 175*a*c\^{}3*s\^{}2\ -\ 1446*a\^{}4*s\ -\ 3036*a\^{}3*c*s\ -\ 1151*a\^{}2*c\^{}2*s\ +\ 438*a*c\^{}3*s\ -\ c\^{}4*s\ -\ 318*a\^{}4\ -\ 795*a\^{}3*c\ -\ 477*a\^{}2*c\^{}2\ +\ 159*a*c\^{}3\ +\ 159*c\^{}4}

\texttt{g\_3}: 136 terms, see \texttt{certificates/G\_3.json} (and \texttt{g\_4}: \texttt{certificates/G\_4.json}). The JSON format is
\texttt{\{"vars":\ {[}"a","c","s"{]},\ "terms":\ {[}{[}e\_a,\ e\_c,\ e\_s,\ "p/q"{]},\ ...{]}\}} and stores \(G_k\) itself (with the constant
and the monomial factor).

\end{document}
